% =====================================================================================================
% PART 3 (frames 10-14): the negation trick and the proof of Theorem 1.1.
% Sources in the paper: Section 1.2 (paragraph "The negation trick"), Section 3 (Figure 3, Figure 4 (a),(b),
% Lemma 3.1, proof of Theorem 1.1, Remarks 3.2 and 3.3).
% Figures: figs/p3-tuples.tex, figs/p3-gadgets-EN.tex, figs/p3-pertype.tex
% =====================================================================================================

% local helpers of part 3 (names are prefixed with pIII to avoid clashes with other parts)
\providecommand{\pIIIbox}[1]{\tikz[baseline=(b.base)]\node[draw=acc,line width=0.8pt,fill=acc!5,
  rounded corners=3pt,inner xsep=9pt,inner ysep=6pt] (b) {#1};}
\providecommand{\pIIInum}[1]{\tikz[baseline=(n.base)]\node[circle,fill=acc,text=white,font=\bfseries\large,
  inner sep=0pt,minimum size=8.5mm] (n) {#1};}
\providecommand{\pIIIto}{\ensuremath{\longrightarrow}}
% "result chip" of frame 13: a fraction of the total weight; #1 = colour, #2 = fraction
\providecommand{\pIIIchip}[2]{\tikz[baseline=(c.base)]\node[fill=#1!7,rounded corners=3pt,inner xsep=4pt,
  inner ysep=3.5pt,align=center] (c) {{\Large\color{#1}$\approx\boldsymbol{#2}$}\\[1.1mm]{\small of the weight}};}
% three small vertices, exactly two of them with the same colour (fills as on frame 9)
\providecommand{\pIIItwoeq}{\tikz[baseline=-0.6ex]{\foreach \i/\c in {0/black!35,1/black!35,2/acc}
  \node[circle,draw,line width=0.6pt,inner sep=0pt,minimum size=9pt,fill=\c] at (0.4*\i,0) {};}}

% ---- Frame 10: The 27 colour patterns of an equation
\begin{frame}{The 27 colour patterns of an equation}
\vspace{0.5mm}
\centering
{\small\color{soft}values $x=x_1x_2x_3\in\Z_3^3$ of the three variables;\quad $s=x_1+x_2+x_3$}\par
\vspace{2.6mm}
% Frame 10: the 27 tuples of Z_3^3 (x written as x_1x_2x_3) in the classes A, D, T_1, T_2.
% Adapted from ../paper/figures/tuples.tex (Figure 3 of the paper): same boxes, same tuples at the same
% positions, same class names/descriptions, same sat_t braces, same three negation arrows, same highlighted
% pair 121 <-> 212.  Changes for the slide: larger fonts and geometry; the two bottom rows (V_E, V_N) are
% left out (their content is on frame 12); the label "x -> -x" is replaced by the example "121 -> 212"
% from the paper's caption (the slide has the line "Negation x -> -x ..." directly under the figure).
% Box geometry in x-units as in the paper: A=[0,1], D=[1.75,3.75], T_1=[4.5,7.5], T_2=[8.25,11.25].
\begin{tikzpicture}[x=1.2cm,y=1cm,
  cls/.style={draw,rounded corners=2.5pt,line width=0.7pt,fill=black!3},
  tup/.style={font=\large,inner sep=0pt},
  hlt/.style={tup,text=acc},
  cname/.style={font=\Large,inner sep=1pt,text height=1.7ex,text depth=0.5ex},
  cdesc/.style={font=\small,text=black!75,inner sep=1pt,text height=1.6ex,text depth=0.5ex},
  neg/.style={<->,>={Stealth[length=5.5pt,width=4.6pt]},acc,line width=0.9pt},
  brc/.style={decorate,decoration={brace,amplitude=5pt},line width=0.6pt},
  slab/.style={above=5pt,font=\normalsize,inner sep=1pt,text height=1.7ex,text depth=0.5ex},
  eg/.style={font=\small,text=acc,inner sep=1pt}]
\def\rh{0.66}\def\bt{2.18}% row height; top of the boxes (3*0.66 + 2*0.1)
% ---------------------------------------------------------------- boxes
\foreach \l/\r in {0/1, 1.75/3.75, 4.5/7.5, 8.25/11.25}
  \draw[cls] (\l,0) rectangle (\r,\bt);
% ---------------------------------------------------------------- tuples
\foreach \i/\a/\d/\e/\p/\q/\r/\pp/\qq/\rr in {
    0/000/012/021/001/010/100/002/020/200,
    1/111/102/201/112/{}/211/221/{}/122,
    2/222/120/210/022/202/220/011/101/110}{
  \pgfmathsetmacro{\yy}{\bt-0.1-\i*\rh-0.5*\rh}
  \node[tup] at (0.5,\yy)   {$\a$};
  \node[tup] at (2.25,\yy)  {$\d$};
  \node[tup] at (3.25,\yy)  {$\e$};
  \node[tup] at (5,\yy)     {$\p$};
  \node[tup] at (6,\yy)     {$\q$};
  \node[tup] at (7,\yy)     {$\r$};
  \node[tup] at (8.75,\yy)  {$\pp$};
  \node[tup] at (9.75,\yy)  {$\qq$};
  \node[tup] at (10.75,\yy) {$\rr$};}
% the T_2 entry at each position is the negation of the T_1 entry there; one pair is highlighted
\node[hlt] at (6,\bt-0.1-1.5*\rh)    {$\mathbf{121}$};
\node[hlt] at (9.75,\bt-0.1-1.5*\rh) {$\mathbf{212}$};
% ---------------------------------------------------------------- headers
\foreach \c/\n/\d in {0.5/{\mathcal A}/all equal, 2.75/{\mathcal D}/all distinct,
                      6/{\mathcal T_1}/{two equal, $s=1$}, 9.75/{\mathcal T_2}/{two equal, $s=2$}}{
  \node[cname] at (\c,\bt+0.86) {$\n$};
  \node[cdesc] at (\c,\bt+0.32) {\d};}
\draw[brc] (0,\bt+1.26)    -- node[slab] {$\sat_0$\ \ ($s(x)=0$)} (3.75,\bt+1.26);
\draw[brc] (4.5,\bt+1.26)  -- node[slab] {$\sat_1$} (7.5,\bt+1.26);
\draw[brc] (8.25,\bt+1.26) -- node[slab] {$\sat_2$} (11.25,\bt+1.26);
% ---------------------------------------------------------------- negation
\draw[neg] (0.3,-0.09)  to[out=-60,in=-120,looseness=1.6] (0.7,-0.09);
\draw[neg] (2.45,-0.09) to[out=-50,in=-130,looseness=1.3] (3.05,-0.09);
\draw[neg] (6,-0.09)    to[out=-14,in=-166] (9.75,-0.09);
\node[eg] at (7.875,-0.66) {e.g.\ $121\mapsto212$};
\end{tikzpicture}
\par
\vspace{2.6mm}
\hl{Negation} $x\mapsto -x$ (swap the values $1$ and $2$):\qquad
$\mathcal A\to\mathcal A$,\quad $\mathcal D\to\mathcal D$,\quad
{\color{acc}$\boldsymbol{\mathcal T_1\leftrightarrow\mathcal T_2}$}
\end{frame}

% ---- Frame 11: The negation trick
\begin{frame}{The negation trick}
\vspace{-0.5mm}
\begin{block}{Observation}
If $X$ is an assignment, so is $-X$, and\\[1.2mm]
\centerline{$-X$ satisfies\ \ $x_1+x_2+x_3=t$
\quad$\Longleftrightarrow$\quad
$X$ satisfies\ \ $x_1+x_2+x_3=-t$}
\end{block}
\pause
\vspace{1.2mm}
\bad{\textbf{NO case:}}\ \ every assignment $X$ satisfies a fraction\ \ $\val(X)\le\frac13+\eps$\ \ of the equations.\par
\vspace{0.6mm}
Random equation:\ \ $t$ = its type,\ \ $s$ = the sum of its variables under $X$.
\par\vspace{2mm}
\centerline{$\Pr[\,s=t\,]\;+\;\Pr[\,s=-t\,]\;\;=\;\;\val(X)+\val(-X)\;\;\le\;\;\tfrac23+2\eps$}
\pause
\vspace{2.6mm}
{\color{acc!35}\rule{\linewidth}{0.5pt}}\par
\vspace{1.4mm}
\textbf{Types $t=1$ and $2$:}
\par\vspace{1.6mm}
\centerline{``$s=t$ or $s=-t$''
\ \ $\Longleftrightarrow$\ \ $s\neq0$
\ \ $\Longleftrightarrow$\ \ \textbf{exactly two values are equal}\ \ \pIIItwoeq}
\par\vspace{2mm}
\centerline{\hl{invariant under renaming the colours\ \ \pIIIto\ \ no reference vertices needed}}
\end{frame}

% ---- Frame 12: Two gadgets
\begin{frame}{Two gadgets}
\centering
\vspace{0.5mm}
% Frame 12: the gadgets E (left) and N (right) = panels (a) and (b) of the paper's Figure 4
% (../paper/figures/gadgets.tex), with EXACTLY the same edges and weights; only the geometry and the
% label sizes are enlarged for the slide, the panel captions are replaced by slide text (name, type and
% source at the top left of each gadget, two caption lines under it), and a two-entry legend is added.
% Every edge is drawn explicitly, one \draw per edge; the edge list extracted from these lines was compared
% by a script with panels (a),(b) of the paper's figure and with the definitions in Section 3 of the paper.
% Edge styles: e1, e1p = weight 1;  e3 = weight 3.
\begin{tikzpicture}[x=1cm,y=1cm,
  prim/.style={circle,draw,line width=0.5pt,fill=black!15,inner sep=0pt,minimum size=14.5pt,font=\small},
  aux/.style={circle,draw,line width=0.5pt,fill=white,inner sep=0pt,minimum size=14pt,font=\footnotesize},
  e1/.style={line width=0.55pt},
  e1p/.style={line width=0.55pt,black!60},
  e3/.style={line width=1.7pt},
  gname/.style={font=\LARGE,text=acc,inner sep=0pt,anchor=base west},
  gsub/.style={font=\small,inner sep=0pt,anchor=base west},
  capA/.style={font=\small,inner sep=0pt,anchor=base,text=black!80},
  capB/.style={font=\normalsize,inner sep=0pt,anchor=base},
  leg/.style={font=\footnotesize,inner sep=0pt,anchor=west,text=black!80}]
% ================================================================== (a) E
\begin{scope}[shift={(0,0)}]
  \node[aux] (aA1) at (90:1.9) {$a_1$};
  \node[aux] (aA4) at (30:1.9) {$a_4$};
  \node[aux] (aA5) at (-30:1.9) {$a_5$};
  \node[aux] (aA2) at (-90:1.9) {$a_2$};
  \node[aux] (aA3) at (-150:1.9) {$a_3$};
  \node[aux] (aA6) at (150:1.9) {$a_6$};
  \node[prim] (pA1) at (90:0.855) {$p_1$};
  \node[prim] (pA2) at (-30:0.855) {$p_2$};
  \node[prim] (pA3) at (210:0.855) {$p_3$};
  % the 6-cycle a1 a4 a5 a2 a3 a6
  \draw[e1] (aA1) -- (aA4);
  \draw[e1] (aA4) -- (aA5);
  \draw[e1] (aA5) -- (aA2);
  \draw[e1] (aA2) -- (aA3);
  \draw[e1] (aA3) -- (aA6);
  \draw[e1] (aA6) -- (aA1);
  % each primary is adjacent to four auxiliary vertices
  \draw[e1p] (pA1) -- (aA1);
  \draw[e1p] (pA1) -- (aA2);
  \draw[e1p] (pA1) -- (aA3);
  \draw[e1p] (pA1) -- (aA4);
  \draw[e1p] (pA2) -- (aA3);
  \draw[e1p] (pA2) -- (aA4);
  \draw[e1p] (pA2) -- (aA5);
  \draw[e1p] (pA2) -- (aA6);
  \draw[e1p] (pA3) -- (aA1);
  \draw[e1p] (pA3) -- (aA2);
  \draw[e1p] (pA3) -- (aA5);
  \draw[e1p] (pA3) -- (aA6);
  % slide text
  \node[gname] at (-3.5,1.52) {$E$};
  \node[gsub]  at (-3.5,0.96) {type $0$};
  \node[gsub]  at (-3.5,0.48) {\src{NRT25}};
  \node[capA] at (0,-2.82) {$6$ auxiliary vertices, $18$ unit edges};
  \node[capB] at (0,-3.44) {max cut $=$ \hl{$18$ if $s=0$}, else $16$};
\end{scope}
% ================================================================== (b) N
\begin{scope}[shift={(7.4,-0.525)}]
  \node[prim] (pB1) at (90:2.0) {$p_1$};
  \node[prim] (pB2) at (-30:2.0) {$p_2$};
  \node[prim] (pB3) at (210:2.0) {$p_3$};
  \node[aux] (qB) at (0,0) {$q$};
  \draw[e1] (pB1) -- (pB2);
  \draw[e1] (pB2) -- (pB3);
  \draw[e1] (pB3) -- (pB1);
  \draw[e3] (qB) -- (pB1);
  \draw[e3] (qB) -- (pB2);
  \draw[e3] (qB) -- (pB3);
  % slide text
  \node[gname] at (-3.4,2.045) {$N$};
  \node[gsub]  at (-3.4,1.485) {types $1$ and $2$};
  \node[gsub,text=acc]  at (-3.4,1.005) {this paper};
  \node[capA] at (0,-2.295) {triangle $+$ vertex $q$;\ \ thick edges: weight $3$};
  \node[capB] at (0,-2.915) {max cut $=$ \hl{$11$ if $s\neq0$}, else $9$};
\end{scope}
% ================================================================== legend (between the gadgets)
\node[prim,minimum size=9pt] at (2.55,-0.78) {};
\node[leg] at (2.85,-0.78) {variable vertices};
\node[aux,minimum size=9pt]  at (2.55,-1.28) {};
\node[leg] at (2.85,-1.28) {auxiliary vertices};
% pointer to the paper, between the two "max cut" lines (both values are Lemma 3.1)
\node[font=\footnotesize,inner sep=0pt,anchor=base] at (3.7,-3.44) {\paperref{Lemma 3.1}};
\end{tikzpicture}
\par
\vspace{3.4mm}
{\small\color{soft}$s=x_1+x_2+x_3$:\quad
$s=0$: all equal or all distinct;\quad $s\neq0$: exactly two equal}
\end{frame}

% ---- Frame 13: Proof of Theorem 1.1 on one slide
\begin{frame}{Proof of Theorem 1.1 on one slide}
\vspace{-1mm}
\textbf{Gadgets:}\ \ \hl{$2E$} for type $0$,\ \ \hl{$N$} for types $1$ and $2$
\hfill{\small\color{soft}total edge weight: $\tfrac13(36+12+12)=20$}
\par\vspace{0.6mm}
\centerline{\begin{tikzpicture}[baseline=(b.base)]
\node[draw=acc,line width=0.8pt,fill=acc!5,rounded corners=3pt,inner xsep=9pt,inner ysep=5.5pt] (b)
  {$V_t(x)\;=\;\beta_t+2\,\big([\,s=t\,]+[\,s=-t\,]\big),\qquad \beta_0=32,\ \ \beta_1=\beta_2=9$};
\node[anchor=south west,font=\footnotesize,text=soft,inner sep=1.5pt] at ([xshift=1.5pt]b.north west)
  {$V_t(x)$: max cut of the type-$t$ gadget on pattern $x$};
\end{tikzpicture}}
\pause
\par\vspace{2.2mm}
\begin{minipage}[c]{116mm}\raggedright
\hl{YES case:}\ \ an assignment that satisfies $\ge1-\eps$ gives a colouring with\par
\vspace{1.4mm}
\hspace*{10mm}$\cut\;\ge\;\tfrac13(36+11+11)-O(\eps)\;=\;\tfrac{58}{3}-O(\eps)$
\end{minipage}\hfill
\begin{minipage}[c]{25mm}\centering\pIIIchip{acc}{\tfrac{29}{30}}\end{minipage}
\pause
\par\vspace{2.2mm}
\begin{minipage}[c]{116mm}\raggedright
\bad{\textbf{NO case:}}\ \ any colouring gives an assignment $\sigma$, and\par
\vspace{1.4mm}
\hspace*{10mm}$\cut\;\le\;\tfrac{50}{3}+2\,\big(\val(\sigma)+\val(-\sigma)\big)
\;\le\;\tfrac{50}{3}+\tfrac43+4\eps\;=\;18+4\eps$
\end{minipage}\hfill
\begin{minipage}[c]{25mm}\centering\pIIIchip{warn}{\tfrac{9}{10}}\end{minipage}
\pause
\par\vspace{2mm}
\centerline{\pIIIbox{\large\color{acc}$\dfrac{9/10}{29/30}\;=\;\boldsymbol{\dfrac{27}{29}}$}}
\end{frame}

% ---- Frame 14: Two remarks
\begin{frame}{Two remarks}
\vspace{1mm}
\pIIInum{1}\hspace{3.5mm}%
\begin{minipage}[t]{128mm}
{\large\textbf{Judged gadget by gadget, $N$ is useless}}\hfill\paperref{Remark 3.3}\par
\vspace{2.2mm}
\begin{minipage}[c]{70mm}\raggedright
As a type-$1$ gadget, it cannot tell $s=1$ from $s=2$.\par
\vspace{2mm}
\hl{The gain comes from the analysis}: pairing $\sigma$ with $-\sigma$.
\end{minipage}\hfill
\begin{minipage}[c]{55mm}\centering
% Frame 14, remark (1): the row "V_N" of the paper's Figure 3 (../paper/figures/tuples.tex), values of
% Lemma 3.1: V_N = 9 on A and D, V_N = 11 on T_1 and on T_2 (larger value shaded, as in the paper).
% Cell geometry in x-units as in the paper: A=[0,1], D=[1.75,3.75], T_1=[4.5,7.5], T_2=[8.25,11.25].
\begin{tikzpicture}[x=0.42cm,y=1cm,
  cell/.style={draw=black!45,line width=0.5pt},
  hi/.style={cell,fill=black!12},
  val/.style={font=\normalsize},
  cname/.style={font=\normalsize,inner sep=1pt,text height=1.6ex,text depth=0.5ex},
  rlab/.style={font=\normalsize,anchor=east,inner sep=1pt},
  note/.style={font=\footnotesize,inner sep=1pt,text height=1.6ex,text depth=0.5ex},
  brc/.style={decorate,decoration={brace,amplitude=4pt,mirror},line width=0.5pt,acc}]
\foreach \c/\n in {0.5/{\mathcal A}, 2.75/{\mathcal D}, 6/{\mathcal T_1}, 9.75/{\mathcal T_2}}
  \node[cname] at (\c,0.88) {$\n$};
\node[rlab] at (-0.35,0.3) {$V_N$};
\foreach \l/\r/\v/\s in {0/1/9/cell, 1.75/3.75/9/cell, 4.5/7.5/11/hi, 8.25/11.25/11/hi}
  \draw[\s] (\l,0.6) rectangle node[val] {$\v$} (\r,0);
\node[note] at (6,-0.27)    {$s=1$};
\node[note] at (9.75,-0.27) {$s=2$};
\draw[brc] (4.5,-0.52) -- node[below=4.5pt,font=\footnotesize,text=acc,inner sep=1pt] {same value} (11.25,-0.52);
\end{tikzpicture}

\end{minipage}
\end{minipage}

\vspace{6mm}
\pIIInum{2}\hspace{3.5mm}%
\begin{minipage}[t]{128mm}
{\large\textbf{The argument is robust}}\hfill\paperref{Remark 3.2}\par
\vspace{2.2mm}
\begin{minipage}[c]{70mm}\raggedright
It needs neither role separation nor, essentially, balance.
\end{minipage}\hfill
\begin{minipage}[c]{55mm}\centering
{\small if type $0$ has weight $\lambda_0$, the ratio is}\par
\vspace{2.2mm}
$\dfrac{23\lambda_0+\frac{31}{3}}{25\lambda_0+11}\;\le\;\dfrac{31}{33}\;<\;\dfrac{16}{17}$
\end{minipage}
\end{minipage}
\end{frame}
